% Generated draft: requires mathematical and source review.
\documentclass[11pt]{article}
\usepackage[margin=25mm]{geometry}
\usepackage{fontspec}
\setmainfont{Latin Modern Roman}
\usepackage{amsmath,amssymb,mathtools,unicode-math}
\setmathfont{Latin Modern Math}
\usepackage{enumitem,needspace,xurl,hyperref}
\setlength{\parindent}{0pt}
\setlength{\parskip}{5pt}
\setlength{\emergencystretch}{3em}
\newcommand{\N}{\mathbb{N}}
\newcommand{\Z}{\mathbb{Z}}
\newcommand{\Q}{\mathbb{Q}}
\newcommand{\R}{\mathbb{R}}
\newcommand{\C}{\mathbb{C}}
\newcommand{\F}{\mathbb{F}}
\newcommand{\A}{\mathbb{A}}
\newcommand{\PP}{\mathbb{P}}
\newcommand{\E}{\mathbb{E}}
\newcommand{\eps}{\varepsilon}
\newcommand{\dd}{\,\mathrm{d}}
\begin{document}
\section*{UnsolvedMath ID: 800378}
\textbf{Source problem code:} catalogue entry
\subsection{Definitions and mathematical statement}
Let \(S\) be an arbitrary set, let \(T\subseteq S\) be fixed, and let \(f:S\to S\). Write \(\N_0=\{0,1,2,\ldots\}\), define \(f^0=\operatorname{id}_S\) and \(f^{n+1}=f\circ f^n\), and, for each \(\alpha\in S\), define the hitting-time set
\[
H_{f,T}(\alpha)=\{n\in\N_0:f^n(\alpha)\in T\}.
\]
The supplied question asks for a classification of maps \(f\) for which every \(H_{f,T}(\alpha)\) has a suitably nice structure, giving finite unions of one-sided arithmetic progressions as an example. The phrase ``nice structure'' is not defined, so this broader request does not specify a unique mathematical property.

Under the positive-step convention, the illustrative classification problem is to characterize all maps \(f:S\to S\) satisfying
\[
\forall\alpha\in S\quad
\exists r_\alpha\in\N_0\quad
\exists (a_{\alpha,j},d_{\alpha,j})_{j=1}^{r_\alpha}
\in\bigl(\N_0\times\{1,2,3,\ldots\}\bigr)^{r_\alpha}
\quad
H_{f,T}(\alpha)
=\bigcup_{j=1}^{r_\alpha}
\{a_{\alpha,j}+k d_{\alpha,j}:k\in\N_0\}.
\]
The union is empty when \(r_\alpha=0\); its parameters may depend on \(\alpha\). The source does not specify whether singleton progressions are permitted. If they are permitted, the displayed condition instead allows \(d_{\alpha,j}=0\), with \(\{a_{\alpha,j}+k\cdot0:k\in\N_0\}=\{a_{\alpha,j}\}\). These conventions give different properties and remain unresolved here. No additional structure on \(S\), restriction on \(T\), or regularity hypothesis on \(f\) is supplied.
\subsection{Short English statement}
For a fixed subset \(T\) of a set \(S\), classify self-maps whose every forward orbit visits \(T\) at times having a specified simple structure. The source suggests finite unions of one-sided arithmetic progressions, but leaves both the broader structural target and the treatment of singleton progressions unspecified.
\subsection{Sources}
\begin{itemize}
\item Local UnsolvedMath registry, supplied complete source record: permanent ID \(800378\), exact source problem code \texttt{catalogue entry}, title ``Open Problems in Arithmetic Dynamics --- Question 1.'' The registry labels the problem open; this is a supplied status, not an independently verified assessment.
\item Benedetto, DeMarco, Ingram, Jones et al., \emph{Current Trends and Open Problems in Arithmetic Dynamics} (2018), Question 1. Supplied source URL: \url{https://arxiv.org/abs/1806.04980}. The registry attributes the question to these authors and records proposed year 2018.
\item Supplied provenance records access on 2026-10-16 and extraction by \texttt{arxiv\_latex}. These are supplied provenance statements; the source has not been independently accessed or verified for this draft.
\end{itemize}
\end{document}
